\section*{Discussion}\label{sec:discussion}

\subsection*{Genome evolution in large populations}

Evolutionary processes at large demographic scales are expected to leave distinct
signatures in the genome, several of which are observed in our Pacific acorn
barnacle assembly.
For example, large populations are expected to harbor increased levels of segregating
polymorphism, the product of both the reduction in the stochastic loss due to
genetic drift and the increase in the population-scaled mutation rate.
In accordance, we observe a barnacle genome with extreme levels of genetic diversity,
both at the individual and population level (Figure~\ref{fig:sup_kmers};
Table~\ref{tab:sup_popgen_stats}), a substantial proportion of which is located in
coding sequences.
\vspace{0.1cm}

Similarly, large populations are expected to experience an increase in the efficacy of
purifying selection, which is expected to reduce the levels of diversity and increase 
conservation across functional sequences.
We see these patterns reflected in the \bgla{} genome.
Putatively functional sites, particularly coding exons and tRNAs, show both increased
levels of sequence conservation (Figure~\ref{fig:phastcons}) and a reduction in
population-level genetic diversity (Figure~\ref{fig:feature_pi};
Table~\ref{tab:sup_feature_pi}) compared to less constrained features of the genome.
Notably, the rank ordering of the enrichment of conserved elements, with tRNAs and CDSs
as the most enriched, parallels findings in insects \citep{siepel_evolutionarily_2005},
suggesting that the targets of purifying selection are broadly conserved across
arthropods despite the much larger and more repeat-rich genome of \bgla{}.
Increased efficacy of purifying selection is also reflected in the results of the
McDonald-Kreitman test (Figure~\ref{fig:mkTest}), where we observe a relative decrease
in the proportion of non-synonymous polymorphism in a central Oregon \bgla{} population,
suggesting that weakly deleterious variants are more efficiently purged by selection
\citep{li_reverse_2008}.
Positive selection is also expected to increase in efficacy in large populations,
resulting in the fixation of a larger fraction of beneficial non-synonymous mutations.
We observe this pattern as well: in this central Oregon population, the majority
(54 to 61\%, depending on the metric of $\alpha$ used) of non-synonymous substitutions
reached fixation due to positive selection, a proportion higher than that observed in
diverse \textit{Drosophila} populations \citep{haller_asymptoticmk_2017}.
\vspace{0.1cm}

In large enough populations, selection should be efficacious enough that its effects
extend to synonymous sites, where substitutions are normally considered neutral.
This could result in a bias in the usage of synonymous codons, as different codons are
favored over others due to GC biases, mRNA stability, tRNA abundance, and other factors
\citep{galtier_codon_2018,presnyak_codon_2015}.
These effects are observed in the \bgla{} assembly in the form of reduction in the
effective number of codons seen across protein-coding genes.
On average, \bgla{} uses $\approx$ 43 out of 61 available codons, a greater reduction than
that observed in other metazoan genomes (Figure~\ref{fig:codon_enc};
Table~\ref{tab:sup_enc_summary}).
There is moderate evidence that this bias is not the product of compositional neutral
processes (Figure~\ref{fig:sup_enc_gc3s};~Figure~\ref{fig:sup_enc_reg}), suggesting an
increase in the efficacy of natural selection due to the large size of Pacific acorn
barnacle populations.
Moreover, these biases put into question the neutrality of at least some synonymous
substitutions in this system, which could have important ramifications for the study of
molecular evolution; however, the extent of these effects remain presently unknown and
should be the subject of future study.
\vspace{0.1cm}

In contrast, the distribution of transposable and repeat elements in the
\bgla{} assembly was unexpected, as we find a highly repetitive genome
composed of over 60\% repeats (Figure~\ref{fig:genome_pi};
Table~\ref{tab:sup_repeat_summary}).
Previous work has shown that purifying selection plays a key role in the accumulation
and fixation of transposable elements \citep{merel_relaxed_2025,langmuller_genomic_2023,
bourque_ten_2018}, following the expectation that repetitive sequences are, on average,
slightly deleterious.
Given the increased efficacy of selection we observed in the rest of the genome,
our expectation was to find a relatively small proportion of repeats.
However, the high repeat content observed in the genome does not necessarily
reflect the low efficacy of selection in these populations.
Some could be low-frequency, slightly deleterious variants in the process of being purged
from the population.
Other repeats might represent beneficial genetic variation, as transposable elements are
known to play roles in adaptive evolution \citep{schrader_impact_2019}.
For example, repeat elements can act as regulators of gene expression
\citep{trizzino_transposable_2017,rech_population-scale_2022}, and as such may provide
the regulatory plasticity required for survival in rapidly changing
intertidal environments.
Regardless of their role, disentangling the population-wide repeat landscape and its
fitness effects is likely impossible with a single haploid reference assembly.
Thus, this observation drives the need for larger comparative genomic studies across
barnacle populations before we can fully understand the dynamics of genomic repeats
and transposable elements in this and other large populations.

\subsection*{Molecular evolution in large populations: reconciling \dnds{} and the
McDonald-Kreitman test}\label{sec:discussion_dnds_mkt}

One of the more striking patterns in our results is the apparent discordance between
our two measures of selection on protein-coding genes.
The average pairwise \dnds{} between \bgla{} and \textit{B. crenatus} is 0.206,
indicating that the majority of coding sites are under purifying selection.
Yet the McDonald-Kreitman test reveals a seemingly opposite pattern: $\approx61\%$
of non-synonymous substitutions fixed between these species were driven by
positive selection ($\alpha_\text{asym}$ = 0.610).
At first glance, these results appear contradictory: how can coding sequences be
under strong purifying selection while simultaneously showing pervasive adaptive
evolution?
\vspace{0.1cm}

The resolution lies in the fact that \dnds{} and the McDonald-Kreitman test measure
fundamentally different quantities. The \dnds{} ratio averages across all
non-synonymous sites in a gene, the vast majority of which are under functional
constraint, resulting in a ratio well below one. In contrast, the McDonald-Kreitman
test partitions changes into polymorphism and divergence, using synonymous changes
as a neutral baseline to ask whether there is an \textit{excess} of non-synonymous
fixations relative to the neutral expectation. This test is therefore sensitive to
the subset of non-synonymous mutations that escape purifying selection and reach
fixation, even when most non-synonymous sites in a gene are constrained.
Moreover, pairwise \dnds{} is calculated from a single sequence per species and
therefore incorporates segregating polymorphisms into its divergence estimate. In
species with many weakly deleterious non-synonymous variants segregating at
appreciable frequencies, this can inflate the non-synonymous divergence and distort
estimates of selection \citep{li_reverse_2008}. By explicitly separating
polymorphism from fixation, the McDonald-Kreitman test avoids this conflation.
\vspace{0.1cm}

\citet{li_reverse_2008} compared the performance of \dnds{} and the McDonald-Kreitman
test across human, \textit{Drosophila}, and yeast genomes and found that the two
approaches identify distinct but correlated sets of genes as targets of selection.
Notably, the McDonald-Kreitman test was more powerful than \dnds{} only in
\textit{Drosophila}, the species with the highest polymorphism, where it
identified nearly three times as many genes under positive selection. In the
lower-diversity human and yeast datasets, the pattern was reversed. This finding
suggests that the statistical power of the McDonald-Kreitman test is greatest in
high-polymorphism species, making it particularly well suited for systems like
barnacles where segregating variation is abundant.
\vspace{0.1cm}

The pattern we observe in \bgla{} is thus consistent with theoretical expectations
for large populations: strong purifying selection efficiently removes the majority
of deleterious non-synonymous mutations, keeping \dnds{} low, while at the same time
the increased efficacy of positive selection drives a disproportionate fraction of
the remaining non-synonymous substitutions to fixation, producing a high
$\alpha$. This coupling of strong purifying and positive selection is a hallmark
of evolution in large populations and has been observed in \textit{Drosophila}, where
$\alpha_\text{asym}$ is approximately 0.57 \citep{messer_frequent_2013}.
Our estimate of $\alpha_\text{asym}$ of 0.610 in \bgla{} is comparable to
or slightly higher than that in \textit{Drosophila}, consistent with the expectation
that the even larger effective population size in barnacles further increases the
efficacy of both modes of selection.
\vspace{0.1cm}

\subsection*{Barnacles as models for evolution at biological extremes}

Our work shows that a Pacific acorn barnacle population from central Oregon exhibits
genetic diversity of over 5\% (mean genome-wide $\pi$ of 0.0506).\comment{[Old pixy value; pixy dropped invariant sites with QUAL $<$ 30, so $\pi$ and $\theta_W$ are about 2$\times$ too high. Corrected genome-wide $\pi$ = 0.023 (\texttt{pg\_gpu\_stats}). Update.]}
These levels of genetic diversity are higher than those observed for other model
systems known for their relatively large population sizes.
For example, the genome-wide $\pi$ we observe in this population is approximately
seven-fold\comment{[Old pixy value; pixy dropped invariant sites with QUAL $<$ 30, so $\pi$ and $\theta_W$ are about 2$\times$ too high. With the corrected $\pi$ this ratio is about 3-fold. Update.]} higher than that observed in sub-Saharan populations of
\textit{D. melanogaster} (Figure~\ref{fig:sup_pi_comp})
\citep{nunez_footprints_2025,li_comparative_2025,pool_population_2012,
langley_genomic_2012}.
Indeed, even the $\pi$ observed at 0-fold degenerate sites, the least diverse and
among the most constrained sites in the genome, exceeds the genome-wide diversity
observed in these \textit{Drosophila} populations.\comment{[Old pixy value; pixy dropped invariant sites with QUAL $<$ 30, so $\pi$ and $\theta_W$ are about 2$\times$ too high. The corrected 0-fold $\pi$ (0.0067) no longer exceeds the \textit{Drosophila} genome-wide mean (0.0076, Figure~\ref{fig:sup_pi_comp}); remove or reword this sentence. Update.]}
While elevated, this level of diversity is not unique to barnacles.
Similar magnitudes of genome-wide $\pi$ ($\approx10^{-2}$) have been observed
in other metazoan taxa, including corals \citep{meziere_corals_2025}, bivalves
\citep{liu_bivalves_2025}, mosquitoes \citep{kent_aedes_2025,ag1000g}, nematodes
\citep{teterina_nematodes_2023}, and lancelets \citep{braso-vives_highly_2026}.
The bivalve values, however, were computed from variant sites only, after the removal
of sites with a minor allele frequency below 5\% \citep{liu_bivalves_2025}; they are
therefore not per-site diversities and are not directly comparable to ours.
Moreover, as with other marine invertebrates, barnacle life history traits are likely
direct drivers of the elevated polymorphism observed, as factors such as large census
sizes, small propagule size, large brood size, long larval dispersal, and broad geographic
distribution have been positively correlated with high genetic diversity in these
organisms \citep{leffler_revisiting_2012,romiguier_comparative_2014}.
\vspace{0.1cm}

The magnitude of diversity observed in \bgla{} solidly places them among
hyperpolymorphic populations, those exhibiting diversity above 5\%\comment{[Old pixy value; pixy dropped invariant sites with QUAL $<$ 30, so $\pi$ and $\theta_W$ are about 2$\times$ too high. With corrected values, 4-fold $\pi$ = 4.2\% and genome-wide $\pi$ = 2.3\%; reword. Update.]}
\citep{cutter_molecular_2013}. This adds barnacles to a growing list of emerging
hyperpolymorphic model systems, which are key to studying the limits of biological
variation and molecular evolution in natural populations. Moreover, our observations
of extreme genetic diversity also have important implications for the study of species
boundaries. For example, recent work by~\citet{hart_genomic_2025} suggests that levels
of polymorphism above 1\% neutral diversity in eukaryotes likely serve as biological limits
for the definition of species. Barnacles and other hyperpolymorphic populations put these
putative cutoffs into question. We observe levels of genome-wide diversity of over 5\%
in just a very small portion of the biological range of the Pacific acorn barnacle.
Nucleotide diversity is even higher ($\approx8\%$\comment{[Old pixy value; pixy dropped invariant sites with QUAL $<$ 30, so $\pi$ and $\theta_W$ are about 2$\times$ too high. Corrected 4-fold $\pi$ = 4.2\%. Update.]}) at 4-fold degenerate, and putatively
neutral, sites (Figure~\ref{fig:feature_pi}; Table~\ref{tab:sup_feature_pi}).
Given the species' broader geographical and ecological distribution, genetic
diversity is likely to increase when looking across a larger spatial sampling.
This observation reiterates the importance of future studies of genetic diversity
of this species over its whole geographic range, providing some insights on
the mechanisms that maintain population and species boundaries in the presence
of extreme genetic diversity.
\vspace{0.1cm}

Beyond species boundaries, the extreme diversity observed in barnacles also speaks to
a long-standing question in population genetics: the relationship between effective
and census population sizes. From our results on diversity, we
observe that this barnacle population has an effective size of approximately $10^{6}$.\comment{[Old pixy value; pixy dropped invariant sites with QUAL $<$ 30, so $\pi$ and $\theta_W$ are about 2$\times$ too high. Recompute with the corrected $\theta_W$; see also the sea star rate below. Update.]}
However, these populations also show densities of $10^{4}$ individuals per square meter
\citep{menge_recruitment_2000}. Based on this estimate, the census size of the Pacific
acorn barnacle population in central Oregon should be, conservatively, in the billions of
individuals. While the general disparity between the effective and census size estimates
has been highlighted for decades, i.e., Lewontin's Paradox \citep{lewontin_genetic_1974},
the mechanisms underlying the differences have been long debated
\citep{gillespie_role_1999,gillespie_genetic_2000,smith_hitch-hiking_1974,
kimura_neutral_1983} and are still the subjects of study \citep{buffalo_quantifying_2021}.
Our work on the Pacific acorn barnacle provides a unique opportunity to empirically
study this disparity and the biological processes that maintain it.
\vspace{0.1cm}

\subsection*{Barnacles among highly fecund marine invertebrates}

\llabel{marine_sweepstakes}\revpoint{2}{5}With census sizes in the billions and high fecundity,
\bgla{} belongs to the group of marine invertebrates in which the relationship between
census size, effective size, and diversity has been debated for decades.
High fecundity combined with high early mortality can produce a very large variance in
reproductive success, such that a small fraction of adults produces most of the recruits
in each generation.
This sweepstakes reproductive success predicts ratios of effective to census size far
below one and genetic differences among cohorts and sites at small spatial scales
(chaotic genetic patchiness) \citep{hedgecock_does_1994}, and, when the variance is
extreme, genealogies with multiple mergers rather than the Kingman coalescent
\citep{eldon_coalescent_2006}.
Multiple-merger genealogies give an excess of rare variants and can bias estimates of
effective size and demographic history made under the Kingman coalescent.
The negative genome-wide Tajima's $D$ we observe is consistent with population growth,
with linked selection, and with sweepstakes reproduction, and our sample of six
individuals cannot separate these processes.
Our estimates of \Ne{} should therefore be read as the effective size of an equivalent
Kingman population.
\vspace{0.1cm}

\llabel{marine_oysters}Bivalves provide the closest precedent for the patterns in coding sequences that we
report.
In the Pacific oyster, \citet{sauvage_single_2007} found among the highest densities of
synonymous polymorphism then known in animals and studied their relation to codon
usage.
In the European flat oyster, \citet{harrang_high_2013} found a high load of non-neutral
amino-acid polymorphisms: the ratio of non-synonymous to silent diversity was 0.38, at
a synonymous diversity of only 0.006, and the authors related this load to sweepstakes
reproduction.
\bgla{} does not follow this pattern.
Its synonymous diversity is about seven times higher ($\pi_S = 0.042$) and its
$\pi_N/\pi_S$ is about one third as high (0.11).\comment{[Corrected values, from the MK
alignments with all invariant sites counted (revision\_analysis
\texttt{results/comparative}); they replace the pixy values used elsewhere in the
manuscript.]}
This value is inside the range predicted from the synonymous diversity of 86 animal
species, and inside the range of the marine invertebrates among them
\citep{romiguier_comparative_2014,galtier_adaptive_2016}
(Figure~\ref{fig:sup_comparative}).
A high load of amino-acid polymorphism is thus not a general property of highly fecund
marine invertebrates, and in \bgla{} purifying selection on amino-acid variation is as
efficient as its synonymous diversity predicts.
\vspace{0.1cm}

\llabel{marine_sv}Marine invertebrate genomes also carry extensive structural variation.
Gene presence-absence variation is widespread in the Mediterranean mussel
\citep{gerdol_massive_2020}; duplications cover 16.5\% of the eastern oyster genome
across 60 individuals \citep{modak_extensive_2021}; sequences present in only one
haplotype make up about 69\% of each haplotype of a \textit{Crassostrea sikamea}
genome \citep{liu_bivalves_2025}; inversion polymorphisms segregate in \textit{Ciona}
\citep{satou_chromosomal_2021}; and copy number variants make up about one fifth of
sequenced loci in the American lobster, where they, and not SNPs, are associated with
sea temperature \citep{dorant_copy_2020}.
The \bgla{} genome shows signs of the same complexity: a $k$-mer heterozygosity above
3\% (Figure~\ref{fig:sup_kmers}), about 240\,Mbp of haplotig sequence removed during
assembly (Table~\ref{tab:sup_asm_stats}), and a repeat content above 60\%.
Our SNP-based analyses do not capture this variation.
Population samples sequenced with long reads will be necessary to measure the total
genetic variation in \bgla{} and to test whether, as in the lobster, structural
variants carry a large part of its adaptive variation.
\vspace{0.1cm}

\llabel{marine_mutation_rate}Direct estimates of the germline mutation rate now exist for highly
fecund marine invertebrates.
In the crown-of-thorns sea star, $\mu = 9.13 \times 10^{-9}$ per site per generation
(95\% CI $6.5 \times 10^{-9}$ to $1.2 \times 10^{-8}$), from 14 parent-offspring trios
\citep{popovic_high_2024}, three times the rate we used here.
With this rate, our genome-wide diversity gives a long-term \Ne{} of about
$6 \times 10^{5}$ (about $1.2 \times 10^{6}$ from 4-fold degenerate sites), which makes
the gap to census size even larger.\comment{[From corrected $\pi$: genome-wide 0.0231,
4-fold 0.0423.]}
In tunicates, the estimates are indirect, and they suggest a per-generation rate no
higher than in vertebrates \citep{tsagkogeorga_population_2012}.
A high mutation rate thus does not appear to explain the high diversity of these
species, and the gap between effective and census size must come from other processes,
such as linked selection or sweepstakes reproduction.
